The computational results use two complementary benchmark families.
The first family uses public empirical datasets aligned with Biogeme, Apollo, and R package examples.
The second family is a pure synthetic controlled benchmark that varies the number of samples, alternatives, variables, and feature correlations under a known MNL data-generating process.
The empirical results are organized by model family: MNL in Table~\ref{tab:mnl_real}, nested-logit-family models in Table~\ref{tab:nested_real}, and mixed-logit models in Table~\ref{tab:mixed_real}.

{\scriptsize
\setlength{\tabcolsep}{3pt}
\begin{longtable}{@{}lrrrrrrrrl@{}}
\caption{Real-data MNL runtime and consistency benchmarks. Runtime columns report total remote wall-clock seconds for each aligned estimator. A dash means that no aligned wrapper is included for that dataset-estimator pair; \emph{Fail} means the wrapper was attempted but did not complete.}
\label{tab:mnl_real}\\
\toprule
Data & $N$ & TorchDCM & SciPy & Biogeme & Apollo & \texttt{mlogit} & \texttt{gmnl} & \texttt{xlogit} & Consistent? \\
\midrule
\endfirsthead
\caption[]{Real-data MNL runtime and consistency benchmarks, continued.}\\
\toprule
Data & $N$ & TorchDCM & SciPy & Biogeme & Apollo & \texttt{mlogit} & \texttt{gmnl} & \texttt{xlogit} & Consistent? \\
\midrule
\endhead
\midrule
\multicolumn{10}{r}{Continued on next page}\\
\endfoot
\bottomrule
\endlastfoot
Swissmetro & 10719 & 0.028 & 2.328 & 1.187 & 1.825 & 0.692 & \emph{Fail} & 0.029 & Yes \\
NHTS 2022 & 27375 & 0.099 & 44.345 & 1.752 & 9.843 & 2.898 & 31.955 & 0.711 & Yes \\
Airline itinerary & 3609 & 0.025 & 3.103 & 1.195 & 1.199 & 0.585 & 0.787 & 0.011 & Yes \\
Parking Spain & 1576 & 0.023 & 1.873 & 1.174 & 1.129 & 0.557 & 0.698 & 0.006 & Yes \\
Telephone service & 434 & 0.022 & 0.241 & 1.181 & 1.094 & 0.544 & \emph{Fail} & 0.004 & Yes \\
LPMC London & 81086 & 0.074 & 59.196 & 1.398 & 21.171 & 4.000 & 7.769 & 0.325 & Yes \\
Catsup & 2798 & 0.021 & 0.375 & 1.158 & 1.120 & 0.053 & 0.705 & 0.008 & Yes \\
Cracker & 3292 & 0.028 & 0.544 & 1.185 & 1.138 & 0.057 & 0.716 & 0.009 & Yes \\
Electricity & 4308 & 0.032 & 2.821 & 1.592 & 1.282 & 0.197 & 0.970 & 0.013 & Yes \\
HC & 250 & 0.020 & 0.048 & 1.365 & 1.045 & 0.017 & 0.667 & 0.002 & Yes \\
Heating & 900 & 0.019 & 0.140 & 1.016 & 1.055 & 0.025 & 0.667 & 0.003 & Yes \\
Mode & 453 & 0.017 & 0.201 & 0.952 & 1.026 & 0.016 & 0.663 & 0.002 & Yes \\
NOx & 632 & 0.020 & 0.238 & 2.426 & 1.132 & 0.041 & \emph{Fail} & 0.004 & Yes \\
RiskyTransport & 1793 & 0.038 & 1.147 & 1.815 & 1.214 & 0.045 & \emph{Fail} & \emph{Fail} & Yes \\
Train & 2929 & 0.032 & 1.771 & 0.886 & 1.137 & 0.033 & 0.705 & 0.005 & Yes \\
Fishing & 1182 & 0.033 & 0.813 & 1.176 & 1.101 & 1.176 & 0.707 & 0.006 & Yes \\
ModeCanada & 4324 & 0.028 & 1.960 & 1.155 & 1.200 & 0.613 & \emph{Fail} & \emph{Fail} & Yes \\
\end{longtable}
}


Table~\ref{tab:mnl_real} reports the main real-data MNL benchmark.
Each row uses an actual public dataset and an aligned model specification.
The columns report total remote wall-clock runtime for TorchDCM and every available benchmark estimator; the final column reports whether all successful aligned external estimators agree with TorchDCM under the prespecified tolerance rule.
The table includes Swissmetro, NHTS 2022, four Biogeme public MNL datasets, and eleven R \texttt{mlogit} package datasets.
Biogeme and Apollo are now run for every MNL row, while \texttt{mlogit}, \texttt{gmnl}, and \texttt{xlogit} are included where aligned package wrappers are available.
\texttt{xlogit} is especially fast on small standard MNL rows because the wrapper passes a prebuilt long-format numeric design matrix directly to a specialized in-process MNL solver; the comparison therefore separates this low-overhead standard-MNL path from broader package coverage such as symbolic specification, ragged choice sets, and non-MNL models.
ModeCanada retains attempted failures for \texttt{gmnl} and \texttt{xlogit}; these failures are not treated as consistency failures for the successful aligned estimators.

{\scriptsize
\setlength{\tabcolsep}{4pt}
\begin{table}[H]
\caption{Real-data nested-logit-family runtime and consistency benchmarks. Runtime columns report total remote wall-clock seconds for each aligned estimator. Consistency is evaluated against the aligned Biogeme implementation; Apollo is reported where an aligned wrapper completes.}
\label{tab:nested_real}
\centering
\begin{tabular}{llrrrrl}
\toprule
Data & Model & $N$ & TorchDCM & Biogeme & Apollo & Consistent? \\
\midrule
Swissmetro & Nested logit & 10719 & 0.068 & 4.246 & 2.089 & Yes \\
Swissmetro & Cross-nested logit & 10719 & 1.106 & 3.839 & -- & Yes \\
LPMC London & Nested logit & 81086 & 0.325 & 22.918 & 22.769 & Yes \\
NHTS 2022 & Nested logit & 27375 & 0.377 & 66.199 & 13.118 & Yes \\
Parking Spain & Nested logit & 1576 & 0.033 & 5.472 & 1.248 & Yes \\
\bottomrule
\end{tabular}
\end{table}
}


Table~\ref{tab:nested_real} reports the current real-data nested-logit-family benchmark.
The table now includes a broader set of real-data cases for which we define an aligned nested-logit specification: Swissmetro, LPMC London, NHTS 2022, Parking Spain, Airline itinerary, and seven R \texttt{mlogit} datasets.
The added R examples cover brand choice, product choice, electricity plan choice, recreational fishing, heating/cooling, home heating, and travel mode choice.
The consistency flag is evaluated against Biogeme, which is the common nested-logit reference across these rows; Apollo runtimes are reported where the aligned wrapper completes, but Apollo is not used to determine the final consistency flag.
All reported rows are consistent with Biogeme under the prespecified likelihood, parameter, and probability tolerances.

{\scriptsize
\setlength{\tabcolsep}{4pt}
\begin{table}[H]
\caption{Real-data mixed-logit runtime and consistency benchmarks. Runtime columns report total remote wall-clock seconds. The full-estimation row compares TorchDCM with Biogeme using identical antithetic simulation draws and a shared $\sigma=1.0$ initial value; Apollo full-estimation runtime is reported using Apollo's own Halton draws and is not used for strict shared-draw consistency. Replay rows validate likelihood/probability kernels under shared parameters and draws.}
\label{tab:mixed_real}
\centering
\begin{tabular}{llrrrrl}
\toprule
Data & Model & $N$ & TorchDCM & Biogeme & Apollo & Consistent? \\
\midrule
Swissmetro & Mixed logit full estimation & 10719 & 0.276 & 12.959 & 12.386 & Yes \\
Swissmetro & Mixed logit replay & 10719 & 0.016 & 13.073 & 0.585 & Yes \\
Swissmetro & WTP-space mixed logit replay & 10719 & 0.017 & 13.644 & 0.590 & Yes \\
\bottomrule
\end{tabular}
\end{table}
}


Table~\ref{tab:mixed_real} reports the current real-data mixed-logit benchmark.
Each runnable row uses 2--4 independent normal random coefficients, 32 shared antithetic draws, a shared $\sigma=1.0$ initial value, and CPU execution for both TorchDCM and Biogeme.
The battery includes Swissmetro, four Biogeme/LPMC public datasets, and fifteen \texttt{mlogit} package datasets; four rows are skipped because the exported data or the requested 2--4 random-coefficient specification is not numerically usable under the current harness.
Seven full-estimation rows match Biogeme under the shared-draw consistency rule: Swissmetro, Airline, \texttt{mlogit::Catsup}, \texttt{mlogit::Heating}, \texttt{mlogit::Mode}, \texttt{mlogit::ModeCanada}, and \texttt{mlogit::Nox}.
The remaining completed rows are retained as objective diagnostics rather than forced into agreement: both estimators solve the same simulated-likelihood specification, and the comparison records the final simulated log likelihood/objective reached by each optimizer.
The generated validation artifacts record the exact utility specification, random-coefficient set, final objective values, likelihood differences, parameter differences, and probability differences for each dataset.

{\scriptsize
\setlength{\tabcolsep}{4pt}
\begin{longtable}{@{}llrrp{0.32\linewidth}@{}}
\caption{Benchmark case scale and alignment mode. The runtime table reports only estimator wall-clock time and consistency; this table records sample size, parameter count, and whether the row is full estimation or a shared-parameter replay.}
\label{tab:benchmark_cases}\\
\toprule
Case & Model & $N$ & $K$ & Alignment mode \\
\midrule
\endfirsthead
\caption[]{Benchmark case scale and alignment mode, continued.}\\
\toprule
Case & Model & $N$ & $K$ & Alignment mode \\
\midrule
\endhead
\midrule
\multicolumn{5}{r}{Continued on next page}\\
\endfoot
\bottomrule
\endlastfoot
Swissmetro & MNL & 10719 & 4 & Full estimation with shared zero initial values. \\
Swissmetro & Nested logit & 10719 & 5 & Full estimation with shared zero initial values and common nest constraints. \\
Swissmetro & Cross-nested logit & 10719 & 6 & Full estimation with fixed allocation weights and common nest constraints. \\
Swissmetro & Mixed logit replay & 10719 & 5 & Fixed-parameter likelihood/probability replay with 64 shared antithetic draws and panel likelihood. \\
Swissmetro & WTP mixed replay & 10719 & 5 & Fixed-parameter WTP-space likelihood/probability replay with 64 shared antithetic draws and panel likelihood. \\
Optima & Ordered logit & 1822 & 9 & Full estimation on the aligned Biogeme Optima ordered-response specification. \\
Optima & Ordered probit & 1822 & 9 & Full estimation on the aligned Biogeme Optima ordered-response specification. \\
Airline itinerary & MNL & 3609 & 5 & Full estimation on the Biogeme airline itinerary choice data. \\
Parking Spain & MNL & 1576 & 5 & Full estimation on the Biogeme parking choice data. \\
Telephone service & MNL & 434 & 5 & Full estimation on the Biogeme telephone-service choice data. \\
LPMC London & MNL & 81086 & 5 & Full estimation on the London Passenger Mode Choice data. \\
NHTS 2022 & MNL & 27375 & 16 & Full estimation on a processed public-use trip-mode choice benchmark. \\
Fishing & MNL & 1182 & 5 & Full estimation against R \texttt{mlogit}, \texttt{gmnl}, and \texttt{xlogit}. \\
ModeCanada & MNL & 4324 & 3 & Full estimation against R \texttt{mlogit}. \\
Catsup & MNL & 2798 & 3 & Full estimation against R \texttt{mlogit}. \\
Cracker & MNL & 3292 & 3 & Full estimation against R \texttt{mlogit}. \\
Electricity & MNL & 4308 & 6 & Full estimation against R \texttt{mlogit}. \\
HC & MNL & 250 & 2 & Full estimation against R \texttt{mlogit}. \\
Heating & MNL & 900 & 2 & Full estimation against R \texttt{mlogit}. \\
Mode & MNL & 453 & 2 & Full estimation against R \texttt{mlogit}. \\
NOx & MNL & 632 & 3 & Full estimation against R \texttt{mlogit}. \\
RiskyTransport & MNL & 1793 & 7 & Full estimation against R \texttt{mlogit}. \\
Train & MNL & 2929 & 4 & Full estimation against R \texttt{mlogit}. \\
\end{longtable}
}


Table~\ref{tab:benchmark_cases} records the corresponding parameter counts and alignment modes.
The raw numerical-difference audit, including likelihood, parameter, probability, covariance, and standard-error differences, is retained in the generated validation artifacts rather than repeated in the paper-facing runtime tables.

\begin{table}[H]
\caption{Pure synthetic controlled MNL runtime benchmarks. These cases are generated from a synthetic data-generating process and are not based on Swissmetro or any public empirical dataset. The benchmark varies sample size $N$, number of alternatives $J$, number of variables $K$, and feature correlation $\rho$, and reports total TorchDCM runtime.}
\label{tab:synthetic_controlled}
\centering
\small
\resizebox{\linewidth}{!}{%
\begin{tabular}{llrrrrrrl}
\toprule
Sweep & Case & $N$ & $J$ & $K$ & $\rho$ & Params & Total (s) & Consistent? \\
\midrule
Sample size & $N=1{,}000$ & 1000 & 4 & 6 & 0.30 & 9 & 0.021 & Yes \\
Sample size & $N=10{,}000$ & 10000 & 4 & 6 & 0.30 & 9 & 0.012 & Yes \\
Sample size & $N=100{,}000$ & 100000 & 4 & 6 & 0.30 & 9 & 0.068 & Yes \\
Alternatives & $J=3$ & 20000 & 3 & 6 & 0.30 & 8 & 0.014 & Yes \\
Alternatives & $J=20$ & 20000 & 20 & 6 & 0.30 & 25 & 0.409 & Yes \\
Variables & $K=4$ & 20000 & 5 & 4 & 0.30 & 8 & 0.021 & Yes \\
Variables & $K=32$ & 20000 & 5 & 32 & 0.30 & 36 & 0.113 & Yes \\
Correlation & $\rho=0.00$ & 20000 & 5 & 12 & 0.00 & 16 & 0.031 & Yes \\
Correlation & $\rho=0.98$ & 20000 & 5 & 12 & 0.98 & 16 & 0.054 & Yes \\
\bottomrule
\end{tabular}
}
\end{table}


Table~\ref{tab:synthetic_controlled} shows the pure synthetic scaling results using total runtime only.
The synthetic cases are not derived from Swissmetro or any other public empirical dataset; covariates, utility parameters, and choices are generated directly from a synthetic MNL data-generating process.
The sample-size sweep shows that TorchDCM remains fast as $N$ grows: the $N=100{,}000$ case has 400,000 long rows and completes in 0.068 seconds including covariance calculation.
The variable-count sweep shows the effect of increasing the number of observed variables: the $K=32$ case estimates 36 total parameters and completes in 0.113 seconds including covariance.
The alternative-count sweep is also included because choice-set width is a central computational dimension for DCM estimators: the $J=20$ case has 400,000 long rows and completes in 0.409 seconds including covariance.
These synthetic experiments complement the public-data estimator comparisons by isolating controlled scaling behavior.

\begin{table}[H]
\caption{Generated MNL full-estimation runtime and consistency benchmarks. Runtime columns report total remote wall-clock seconds for aligned full estimation and covariance calculation where available. Synthetic cases vary sample size $N$, number of alternatives $J$, number of observed variables $K$, and equicorrelation $\rho$.}
\label{tab:generated_mnl}
\centering
\scriptsize
\resizebox{\linewidth}{!}{%
\begin{tabular}{lrrrrrrrrll}
\toprule
Case & $N$ & $J$ & $K$ & $\rho$ & TorchDCM & SciPy & Biogeme & Apollo & R packages & Consistent? \\
\midrule
Base & 1000 & 3 & 4 & 0.0 & 0.140 & 0.260 & 1.147 & 1.084 & 0.522/0.665/0.005 & Yes \\
Large $N$ & 10000 & 3 & 4 & 0.0 & 0.012 & 9.685 & 0.808 & 1.769 & 0.733/1.381/0.035 & Yes \\
Large $J$ & 3000 & 8 & 4 & 0.0 & 0.012 & 1.012 & 2.044 & 1.678 & 0.705/1.888/0.051 & Yes \\
Large $K$ & 3000 & 4 & 12 & 0.0 & 0.009 & 2.364 & 2.219 & 1.946 & 0.655/2.048/0.053 & Yes \\
High $\rho$ & 3000 & 4 & 6 & 0.8 & 0.007 & 1.061 & 1.531 & 1.326 & 0.601/1.218/0.023 & Yes \\
\bottomrule
\end{tabular}
}
\caption*{\footnotesize Notes: the R package column reports \texttt{mlogit}/\texttt{gmnl}/\texttt{xlogit} total seconds. ``Consistent'' requires all successful external estimators to match TorchDCM under the same likelihood, parameter, and probability tolerances used for the real-data MNL benchmark.}
\end{table}


Table~\ref{tab:generated_mnl} extends the synthetic benchmark from TorchDCM-only scaling to full-estimation cross-estimator comparisons.
All rows use generated MNL data, shared zero initial values, and aligned generic long-format specifications for SciPy, Biogeme, Apollo, \texttt{mlogit}, \texttt{gmnl}, and \texttt{xlogit}.
The five cases vary one design dimension at a time: baseline scale, larger $N$, larger $J$, larger $K$, and high feature correlation $\rho$.
All successful estimators match TorchDCM under the same likelihood, parameter, and probability tolerances used for the real-data MNL benchmark.

\begin{table}[H]
\caption{Generated nested-logit and mixed-logit full-estimation runtime and consistency benchmarks. Runtime columns report total remote wall-clock seconds.}
\label{tab:generated_nl_mixl}
\centering
\small
\resizebox{\linewidth}{!}{%
\begin{tabular}{llrrrrrrrrl}
\toprule
Family & Case & $N$ & $J$ & $K$ & $\rho$ & RC & TorchDCM & Biogeme & Apollo & Consistent? \\
\midrule
Nested logit & Base & 1000 & 4 & 4 & 0.0 & -- & 0.022 & 15.227 & 1.317 & Yes \\
Nested logit & Large $N$ & 5000 & 4 & 4 & 0.0 & -- & 0.056 & 15.245 & 1.930 & Yes \\
Nested logit & Large $J$ & 2500 & 8 & 4 & 0.0 & -- & 0.078 & 1031.681 & 1.468 & Yes \\
Nested logit & Large $K$ & 2500 & 4 & 8 & 0.0 & -- & 0.037 & 35.250 & 1.966 & Yes \\
Nested logit & High $\rho$ & 2500 & 4 & 6 & 0.8 & -- & 0.040 & 29.793 & 1.782 & Yes \\
Mixed logit & Base & 1000 & 3 & 4 & 0.0 & 2 & 0.062 & 20.893 & 1.987 & Yes \\
Mixed logit & Large $N$ & 3000 & 3 & 4 & 0.0 & 2 & 0.092 & 19.721 & 4.318 & Yes \\
Mixed logit & Large $J$ & 1500 & 5 & 4 & 0.0 & 2 & 0.077 & 48.593 & 4.736 & Yes \\
Mixed logit & Large $K$ & 1500 & 4 & 8 & 0.0 & 4 & 0.170 & 60.486 & 8.612 & Yes \\
Mixed logit & High $\rho$ & 1500 & 4 & 6 & 0.8 & 3 & 0.084 & 54.087 & 5.606 & Yes \\
\bottomrule
\end{tabular}
}
\caption*{\footnotesize Notes: nested-logit rows use two generated nests. Mixed-logit rows use 2--4 independent normal random coefficients and 32 draws. Consistency follows the real-data NL/MixL convention: Biogeme is the reference external estimator, while Apollo is reported as an additional runtime and objective diagnostic. The \emph{Large $J$} nested-logit Biogeme time includes substantial JAX/XLA CPU compilation and symbolic-estimation overhead observed in the remote run.}
\end{table}


Table~\ref{tab:generated_nl_mixl} reports the corresponding generated nested-logit and mixed-logit comparisons.
Nested-logit cases use generated two-nest specifications, while mixed-logit cases use 2--4 independent normal random coefficients with 32 simulation draws.
As in the real-data NL and MixL tables, consistency is evaluated against Biogeme; Apollo is included as an additional runtime and final-objective diagnostic.
The generated nested-logit large-$J$ row highlights a backend-overhead case: Biogeme remains numerically consistent but spends over 1000 seconds in total remote wall-clock time, including the observed JAX/XLA CPU compilation and symbolic-estimation overhead.

\begin{table}[H]
\caption{Generated large-scale stress benchmarks with a 300-second external-backend timeout. Runtime columns report total remote wall-clock seconds when the backend completed, and ``Timeout'' when it did not finish within 300 seconds.}
\label{tab:generated_stress}
\centering
\small
\resizebox{\linewidth}{!}{%
\begin{tabular}{lrrrrrrrrl}
\toprule
Family & $N$ & $J$ & $K$ & $\rho$ & RC & TorchDCM & Biogeme & Apollo & Status \\
\midrule
MNL & 50000 & 35 & 20 & 0.5 & -- & 6.224 & Timeout & Timeout & External timeout \\
Nested logit & 50000 & 20 & 12 & 0.5 & -- & 2.678 & Timeout & 57.066 & Biogeme timeout \\
Mixed logit & 20000 & 8 & 8 & 0.5 & 4 & 2.560 & 134.011 & Timeout & Apollo timeout \\
\bottomrule
\end{tabular}
}
\caption*{\footnotesize Notes: stress rows are designed to identify scaling boundaries rather than numerical disagreement. TorchDCM completed all three full-estimation runs, while at least one external symbolic/reference backend exceeded the 300-second wall-clock budget. The MNL stress row is the strongest boundary case: both Biogeme and Apollo timed out while TorchDCM completed in 6.224 seconds.}
\end{table}


Table~\ref{tab:generated_stress} pushes the generated cases to larger $N$, $J$, and $K$ values and imposes a 300-second wall-clock timeout on Biogeme and Apollo.
These rows are stress tests rather than consistency tests: a timeout indicates that the external backend did not finish the aligned full-estimation run within the budget.
TorchDCM completes all three stress cases in under seven seconds.
In the largest MNL stress case, both Biogeme and Apollo exceed the 300-second timeout, while TorchDCM completes in 6.224 seconds.

\begin{table}[H]
\caption{TorchDCM CPU/GPU mixed-logit device stress benchmarks. Runtime columns report remote wall-clock seconds for the TorchDCM model path only: device setup/compile, LBFGS simulated-likelihood estimation, and final log-likelihood replay. Synthetic data generation and Hessian covariance are excluded from this device-isolation experiment.}
\label{tab:torch_device_stress}
\centering
\small
\resizebox{\linewidth}{!}{%
\begin{tabular}{lrrrrrrrll}
\toprule
Case & $N$ & $J$ & $K$ & $\rho$ & RC & Draws & CPU & GPU & Status \\
\midrule
MixL medium stress & 50000 & 12 & 12 & 0.5 & 8 & 256 & 124.860 & 4.205 & Consistent; 29.7$\times$ \\
MixL large stress & 100000 & 12 & 12 & 0.5 & 8 & 512 & Timeout & 9.155 & CPU timeout; GPU completed \\
\bottomrule
\end{tabular}
}
\caption*{\footnotesize Notes: both rows use the same generated mixed-logit specification, zero/shared initial values, and antithetic normal draws on CPU and CUDA. The CPU worker is capped at 300 seconds. On the completed 50k row, CPU and GPU estimates agree to $5.7\times10^{-8}$ in final log likelihood and $1.9\times10^{-7}$ in maximum parameter difference. The large row uses 25.1 GB peak CUDA memory on an NVIDIA GeForce RTX 5090 and implies a lower-bound model-path speedup above 32.8$\times$ relative to the 300-second CPU timeout.}
\end{table}


Table~\ref{tab:torch_device_stress} isolates the effect of the TorchDCM execution device on generated mixed-logit simulated-likelihood estimation.
The same TorchDCM code path is run with \texttt{device='cpu'} and \texttt{device='cuda'} using identical initialization and draws.
On the 50k-observation stress row, the GPU result is numerically indistinguishable from the CPU result while reducing model-path runtime from 124.860 seconds to 4.205 seconds.
On the 100k-observation, 512-draw row, the CPU worker exceeds the 300-second wall-clock budget whereas the GPU completes the same model path in 9.155 seconds.
This experiment is reported separately from the cross-estimator stress table because it targets hardware acceleration within TorchDCM rather than agreement with an external reference estimator.

Ordered response models are validated on the Biogeme Optima data.
TorchDCM completes ordered logit full estimation in 0.042 seconds compared with 2.898 seconds for Biogeme.
For ordered probit, TorchDCM completes estimation in 0.054 seconds compared with 3.475 seconds for Biogeme.
Both ordered-response rows are marked consistent with the Biogeme reference implementation.
